% !TEX root = ../main.tex
\documentclass[../main.tex]{subfiles}%para poder compilar esta file sola
% Chapter 1
\begin{document}

\chapter{Desarrollo}
\label{chap:desarrollo} 

\section{RTL Design}

\subsection{I2C Controller}

La interfaz de control del sensor MT9V111 incluye un bus I2C, que permite la escritura y lectura de los registros internos del sensor. Dentro de los distintos modos de operación, en este proyecto se implementan las funciones 
de escritura y lectura de datos de 16 bits, que es el tamaño de los registros internos. Además, el sensor permite escritura/lectura iterativa en una sola transmisión y es el modo de comunicación que se usará. 
\cite{mt9v111_datasheet} 

\paragraph{Protocolo I2C MT9V111}

\begin{itemize}
    \item Start: El bit de start se da cuando la señal SDA baja mientras SCLK está en alto (Marcador 3 en \ref{fig:waveforms_16bits_rw_frames}).
    \begin{figure}[htbp] 
    \centering
        %\subfloat[Trama completa]{\includegraphics[width=\textwidth]{gfx/waveforms/16bitswriting/16bitswriting.pdf}\label{fig:waveforms_16bitswriting_a}}
        %\\
        \subfloat[Envío de dirección de esclavo]{\fbox{\includegraphics[width=0.9\textwidth]{gfx/waveforms/16bitswriting/16bitswriting_addr.pdf}\label{fig:waveforms_16bitswriting_addr}}}
        \\
        \subfloat[Envío de dirección del registro \textbf{+ W}]{\fbox{\includegraphics[width=0.9\textwidth]{gfx/waveforms/16bitswriting/16bitswriting_reg_addr.pdf}\label{fig:waveforms_16bitswriting_reg_addr}}}
        \\
        \subfloat[Envío de start repetido y de dirección del resitro \textbf{+ R}]{\fbox{\includegraphics[width=0.9\textwidth]{gfx/waveforms/16bitsreading/16bitsreading_rstart.pdf}\label{fig:waveforms_16bitsreading_rstart}}}
        \\
        \subfloat[Envío de valor del registro: MSB \& LSB ]{\fbox{\includegraphics[width=\textwidth]{gfx/waveforms/16bitswriting/16bitswriting_reg_val.pdf}\label{fig:waveforms_16bitswriting_reg_val}}}
    \caption{Tramas para Escritura y Lectura de 16 bits en el bus I2C del sensor MT9V111.} 
    \label{fig:waveforms_16bits_rw_frames} 
    \end{figure}
\end{itemize}

\begin{figure}[htbp] 
    \centering
    \includegraphics[width=0.6\textwidth]{gfx/waveforms/4phases_sclk.pdf}
    \caption{Cronograma de las 4 fases de SCLK.} 
    \label{fig:waveforms_4phases} 
\end{figure}

\todohint[caption={Figuras diagramas}, inline]{Explicar de forma detallada junto con los diagramas de drawio la secuencia de la FSM del I2C Controller.}

\end{document}
